\insertImage{b67bf4ba-580b-4672-b687-8310d1ca30f5.png}{Aquarelle d'une série de panneaux routiers sur une route sinueuse avec des chemins de terre qui se détachent de la route principale et semblent être des raccourcis.}

\chapter{Les causes des échecs}

Personne n'écrit de communiqué de presse pour annoncer que sa libération a échoué.

C'est pour ça que vous connaissez FAVI, Chrono Flex ou Buurtzorg, et que vous ne connaissez pas les dizaines d'entreprises qui ont tenté la même aventure et qui sont revenues en arrière, discrètement, en espérant que personne ne pose trop de questions. Les conférences sont pleines de réussites. Les échecs, eux, se racontent en off, autour d'un café, quand le dictaphone est éteint.

Ce chapitre, c'est le café sans le dictaphone. On va regarder pourquoi ça échoue — pas pour se moquer, mais parce que je suis convaincu d'une chose : si vous comprenez comment on rate une libération, vous avez fait la moitié du chemin pour la réussir.

\section{Les erreurs de mise en œuvre et de gestion}

L'histoire commence presque toujours pareil. Un dirigeant lit un livre — souvent un très bon livre, d'ailleurs —, sort d'une conférence, et annonce le lundi matin que l'entreprise est désormais libérée. Les managers deviennent des « coachs », l'organigramme part à la poubelle, et tout le monde applaudit. Enfin, presque tout le monde.

Six mois plus tard : plus personne ne sait qui décide de quoi, les anciens managers font le même métier qu'avant sans oser le dire, et les décisions urgentes attendent un consensus qui ne vient jamais.

Le problème, c'est qu'une libération, ça ne se décrète pas. Ça se construit.

Supprimer la hiérarchie, c'est la partie facile — une réunion suffit. La partie difficile, c'est de construire ce qui la remplace : des règles de décision explicites, une information réellement partagée, des mécanismes d'arbitrage quand deux équipes ne sont pas d'accord. La hiérarchie rendait des services invisibles ; tant qu'on ne les a pas identifiés, on ne peut pas les remplacer. Les entreprises qui retirent la structure sans rien installer à la place ne récoltent pas la liberté : elles récoltent le flou\footnote{Laloux, F. (2014). Reinventing Organizations: A Guide to Creating Organizations Inspired by the Next Stage of Human Consciousness. Nelson Parker.}.

Et puis il y a l'erreur miroir, plus subtile : libérer les gens malgré eux. On ne demande pas leur avis aux équipes, puisque c'est pour leur bien. Relisez cette phrase — le contresens tient en douze mots. Une transformation qui repose entièrement sur l'autonomie et l'adhésion ne peut pas être imposée d'en haut par celui qui détient le pouvoir. C'est pourtant, très exactement, ce que font beaucoup de dirigeants « libérateurs ».

Je parle en connaissance de cause. Quand j'ai mis en place l'holacratie dans ma propre société, ma plus grosse erreur n'a été ni technique ni organisationnelle : elle a été pédagogique. Je n'ai pas assez dit, dès le départ, qu'il y aurait des cadres — certains déjà en place, d'autres à construire ensemble — et des contraintes de rentabilité tout à fait ordinaires, sans lesquelles le modèle ne peut tout simplement pas vivre. Quand on annonce la liberté sans annoncer le cadre, chacun dessine dans sa tête la liberté qui l'arrange — et la réalité se charge de la gommer. Si c'était à refaire, je commencerais par là : la pédagogie du cadre, avant la promesse de la liberté.

Enfin, il y a tout ce qui relève du pilotage au long cours. Une entreprise libérée n'est pas un projet qu'on livre, c'est un équilibre qu'on entretient : entre l'autonomie des équipes et la cohérence de l'ensemble, entre la liberté de chacun et les tensions que cette liberté fait remonter. Les entreprises qui échouent ne sont pas forcément celles qui ont mal démarré. Ce sont souvent celles qui ont cru que c'était fini.

\section{Les résistances internes et externes}

Parlons maintenant de ce qu'on n'ose pas dire dans les conférences : tout le monde n'a pas envie d'être libéré.

En interne, la résistance la plus visible vient des managers intermédiaires — et franchement, mettez-vous à leur place. On leur a expliqué pendant quinze ans que leur valeur se mesurait à la taille de leur équipe, et du jour au lendemain on leur annonce que leur poste est « un frein à l'autonomie ». Leur résistance n'est pas de la mauvaise volonté : c'est une réaction parfaitement rationnelle de gens à qui on retire leur métier sans leur en proposer un autre.

La résistance plus silencieuse, elle, vient de salariés qui préfèrent un cadre clair, des responsabilités définies, une frontière nette entre le travail et le reste\footnote{Laloux, F. (2014). Reinventing Organizations: A Guide to Creating Organizations Inspired by the Next Stage of Human Consciousness. Nelson Parker.}. Et c'est leur droit le plus strict. Le piège, pour une entreprise en cours de libération, c'est de traiter ces personnes comme des problèmes à convertir. On obtient alors l'exact inverse du but recherché : une entreprise « libérée » où il est devenu obligatoire d'être autonome, enthousiaste, et volontaire pour tout.

À l'extérieur, la liste des sceptiques est longue : des clients qui veulent « parler au responsable » et ne comprennent pas qu'il n'y en ait pas, des banquiers et des investisseurs que l'absence d'organigramme inquiète, des partenaires qui cherchent leur interlocuteur habituel. Ces résistances-là ne sont pas insurmontables — FAVI a continué à gagner la confiance de clients industriels exigeants pendant toute sa transformation\footnote{Getz, I. (2009). Liberating Leadership: How the Initiative-Freeing Radical Organizational Form Has Been Successfully Adopted. California Management Review, 51(4), 32-58.} — mais elles ne se dissolvent pas toutes seules. Il faut quelqu'un pour les entendre, les prendre au sérieux et y répondre.

Une transformation qui échoue, c'est rarement une transformation attaquée. C'est une transformation qui n'a pas écouté.

\section{Les limites du modèle d'entreprise libérée}

Il faut maintenant que je vous dise quelque chose qui ne me fait pas plaisir, à moi qui défends ce modèle : l'entreprise libérée ne convient pas à toutes les organisations. Et prétendre le contraire est probablement la meilleure façon de fabriquer des échecs.

Il y a d'abord les contextes qui s'y prêtent mal. Dans les secteurs très réglementés ou à fortes contraintes de sécurité, une partie des décisions ne peut tout simplement pas être laissée à l'appréciation de chacun — et ce n'est pas un défaut de confiance, c'est le métier. La taille joue aussi : coordonner cinquante personnes sans hiérarchie est un défi stimulant ; en coordonner cinq mille est un problème d'une toute autre nature, que peu d'organisations ont résolu de façon convaincante\footnote{Getz, I. (2009). Liberating Leadership: How the Initiative-Freeing Radical Organizational Form Has Been Successfully Adopted. California Management Review, 51(4), 32-58.}.

Il y a ensuite les risques inhérents au modèle lui-même, même bien appliqué. En distribuant le pouvoir, on distribue aussi la responsabilité — et tout le monde n'est pas également armé pour la porter. Certains s'en saisissent et s'épanouissent. D'autres la subissent en silence, et c'est comme ça qu'un modèle pensé pour le bien-être fabrique du surmenage. On y reviendra longuement dans la quatrième partie, parce que c'est, à mon sens, l'angle mort le plus grave du discours sur la libération.

Alors, faut-il renoncer ? Non. Mais il faut inverser la question. La bonne question n'est pas « comment appliquer le modèle ? », c'est « qu'est-ce que ce modèle peut apporter à \emph{cette} entreprise, avec \emph{ces} personnes, dans \emph{ce} contexte ? ». Les échecs racontés dans ce chapitre ont presque tous un point commun : quelqu'un a répondu à la première question sans jamais s'être posé la seconde.

Le modèle n'a pas échoué. On lui a demandé quelque chose qu'il n'avait jamais promis.
